% Proof of the arbitrary-field virtual-annulus/tropical-ray comparison.
% Date: 27 August 2026.
% This packet is not integrated into the book and makes no novelty claim.

\begin{proof}[Proof of
\cref{prop:p21p3-gf-virtual-annulus-bridge}]
The residual curve and the branch--valuation correspondence over an
arbitrary complete field are
\cite[\S\S3.3.5.2, 4.2.6, 4.2.9, and 4.2.11.1]
{Ducros24CurveStructure}.  Since \(x\)
is interior, every closed point of \(C_x\) occurs as a branch.  For the branch
\(b\) corresponding to \(p_b\), the refined value group is the lexicographic
extension of \(|\mathcal H(x)^{\times}|\) by the discrete valuation
\(\operatorname{ord}_{p_b}\).  If \(s_b\) is a uniformizer of the discrete
valuation ring \(\mathcal O_{C_x,p_b}\), then
\cite[\S4.2.11.1]{Ducros24CurveStructure} says precisely that the class of
its refined value generates
\[
 |\mathcal H(x)^{\times}|_b/|\mathcal H(x)^{\times}|\simeq\mathbf Z.
\]
Since \(Y\) is normal and \(x\) is nonrigid,
\cite[\S4.2.18]{Ducros24CurveStructure} identifies
\(\mathcal O_{Y,x}\) with \(\kappa(x)\), and
\cite[\S4.2.6]{Ducros24CurveStructure} says that \(\kappa(x)\) is dense in
\(\mathcal H(x)\).  Hence \(s_b\), viewed as a homogeneous residue element
of degree one, has a norm-one lift \(f_b\in\mathcal O_{Y,x}\).  This is also
the finite residue-lifting step made explicit in the proof of
\cite[Lemma~6.14]{GJR26}.  Its value in \(\mathcal H(x)\) is nonzero, so
after shrinking it is an analytic unit.  This proves that
\eqref{eq:p21p3-gf-branch-order-definition} is meaningful.  Replacing a
uniformizer by another one multiplies it by a unit in
\(\mathcal O_{C_x,p_b}\), whose class in the displayed quotient is trivial.
This proves independence and also gives
\eqref{eq:p21p3-gf-branch-order-residue}.

We next prove the exact ray assertion.  By
\cite[Proposition~4.3.9(2) and \S\S4.3.10.2--4.3.10.3]
{Ducros24CurveStructure}, \(b\) has a cofinal family of virtual-annular
sections.  Choose one and call it \(Z\).  Its analytic skeleton and canonical
retraction are compatible with every complete scalar extension by
\cite[\S\S3.6.35.5--3.6.35.10]{Ducros24CurveStructure}.  After a finite
separable extension and then passage to a completed algebraic closure, a
component above \(Z\) is an ordinary annulus whose skeleton maps
homeomorphically to \(S^{\mathrm{an}}(Z)\).

Let \(\beta\) be the branch of this split annulus above \(b\), and choose an
annulus coordinate \(\theta\) so that \(-\log|\theta|\) increases away from
the endpoint over \(x\).
The annular calculation in
\cite[\S\S3.6.22--3.6.24 and \S4.3.6.4]
{Ducros24CurveStructure} shows that, after shortening the annulus once for
the finite family \(g_1,\ldots,g_m\),
\begin{equation}
\label{eq:p21p3-gf-split-annulus-monomials}
                  |g_j|=c_j|\theta|^{n_j}
                  \qquad(1\leq j\leq m)
\end{equation}
for constants \(c_j>0\) and integers \(n_j\).  Naturality of the refined
branch valuation is \cite[\S4.2.9]{Ducros24CurveStructure}.  Consequently,
restriction from \(b\) to \(\beta\) sends the generator defined by \(g_m\) to
a positive integral multiple of the primitive annular generator.  Hence
there is one integer \(e_b>0\), independent of \(j\), such that
\begin{equation}
\label{eq:p21p3-gf-common-expansion-index}
                       n_j=e_b\partial_b(g_j)
                       \qquad(1\leq j\leq m).
\end{equation}
Indeed, this is just the homomorphism between the two rank-one quotient
value groups; \(n_m=e_b\) because \(\partial_b(g_m)=1\).  The same integer
is the tangent-lattice expansion index of this split annular presentation
by \cite[\S3.6.24.2]{Ducros24CurveStructure}.  In the finite-morphism
setting of \cite[Theorem~4.3.13(i)]{Ducros24CurveStructure}, that annular
index is identified with the corresponding residual ramification index.

Put
\[
 t=-\log|g_m|+\log|g_m(x)|.
\]
The orientation and
\cite[\S4.2.11.1]{Ducros24CurveStructure} make \(t\) strictly positive
away from \(x\) on a sufficiently short section.  Equations
\eqref{eq:p21p3-gf-split-annulus-monomials} and
\eqref{eq:p21p3-gf-common-expansion-index}, followed by passage to the
limit at \(x\), give on the skeleton
\begin{equation}
\label{eq:p21p3-gf-coordinate-ray-formula}
 -\log|g_j|+\log|g_j(x)|
       =\partial_b(g_j)t
       \qquad(1\leq j\leq m).
\end{equation}
Every invertible analytic function on a virtual annulus has the same norm
as at its canonical retraction by
\cite[\S3.6.35.8]{Ducros24CurveStructure}.  Thus
\eqref{eq:p21p3-gf-coordinate-ray-formula} holds on the whole shortened
section \(Z_b\), and \(t\) identifies its skeleton with an interval
\((0,\varepsilon_b)\).  This proves
\eqref{eq:p21p3-gf-exact-branch-ray}.  Since the last coordinate of \(v_b\)
is one, \(v_b\) is primitive in the saturated ambient lattice.  The
one-dimensional skeleton attached to the moment, restricted to any compact
subannular band, has the unique integral affine structure for which the
moment maps its faces integral-affinely onto their images by
\cite[\S8.13]{GJR26}; hence this is also the primitive tangent of the
corresponding source one-cell.  This proves
\cref{itm:p21p3-gf-bridge:ray}.

We prove the finite separation assertion.  Let \(p_i=p_{b_i}\).  Weak
approximation on the normal projective curve \(C_x\) gives rational functions
\(s_1,\ldots,s_s\) with
\begin{equation}
\label{eq:p21p3-gf-kronecker-orders}
                \operatorname{ord}_{p_i}(s_j)=\delta_{ij}.
\end{equation}
For completeness, choose a closed point \(q\) away from the \(p_i\).  For
\(n\gg0\), Serre vanishing gives
\[
 H^1\!\left(C_x,
   \mathcal O_{C_x}\!\left(nq-2\sum_i p_i\right)\right)=0.
\]
The divisor exact sequence makes the simultaneous value-and-first-jet map
from \(H^0(C_x,\mathcal O_{C_x}(nq))\) onto
\(\bigoplus_i\mathcal O_{C_x,p_i}/\mathfrak m_{p_i}^2\) surjective.
Prescribing a nonzero first jet and value zero at \(p_j\), and value one at
the other \(p_i\), yields \(s_j\) satisfying
\eqref{eq:p21p3-gf-kronecker-orders}.  This argument is valid for closed
points of arbitrary residue degree.

Lift the \(s_j\) to norm-one analytic units by the density argument above
(equivalently, the lifting step in the proof of
\cite[Lemma~6.14]{GJR26}), and first append those lifts to \(F_0\).  By
\cite[Proposition~4.2.4(ii)]{Ducros24CurveStructure}, every
analytic unit has nonzero branch order in only finitely many branches.
The preliminary outgoing vector in branch \(b_i\) therefore has, in the new
\(s_1,\ldots,s_s\) block, the standard basis vector \(e_i\); in particular it
is primitive, and the selected vectors are pairwise nonparallel.

There can still be an unselected branch whose preliminary vector is a
positive multiple of the vector for some \(b_i\).  There are only finitely
many such branches: any one of them has positive \(s_i\)-order and hence
belongs to the finite exceptional set supplied by
\cite[Proposition~4.2.4(ii)]{Ducros24CurveStructure}.  Call their union
\(Q\).  For every \(c\in Q\), weak approximation as above gives a rational
function \(h_c\) satisfying
\begin{equation}
\label{eq:p21p3-gf-collision-killers}
 \operatorname{ord}_{p_{b_i}}(h_c)=0\quad(1\leq i\leq s),
 \qquad
 \operatorname{ord}_{p_c}(h_c)=1.
\end{equation}
Lift all \(h_c\) and append them as well.  The new coordinates vanish in
branch order on every selected branch, whereas the \(h_c\)-coordinate has
order one on \(c\).  Thus each old collision is destroyed.  A newly active
branch cannot create a new collision: projection to the preliminary
coordinates would already have made it an element of \(Q\).  This proves
\cref{itm:p21p3-gf-bridge:separation}.

Finally, consider the localization in
\cref{itm:p21p3-gf-bridge:localization}.  The proof of
\cite[Theorem~16.8.5]{CLD25} starts from a compact domain \(V\), a moment
whose tropical value at \(x\) avoids the tropicalized boundary, and then sets
\(
\Omega=\mathbf R^M\setminus F_{\mathrm{trop}}(\partial V)
\)
and \(U=F_{\mathrm{trop}}^{-1}(\Omega)\).  First shrink to a common
neighborhood on which all residue lifts are defined.  Then reapply
\cite[Lemma~12.3.6]{CLD25} there to choose \(V\) and an auxiliary
boundary-separating moment, and concatenate that moment with all the
coordinates already constructed.  Boundary avoidance is preserved under
this concatenation, because equality for the concatenated moment would
imply equality for its boundary-separating projection.  The branch sections in
\cite[\S4.3.10.3]{Ducros24CurveStructure} are cofinal, so the sections
\(Z_{b_i}\) may be shortened inside \(U\) and inside the already chosen
source-separation plateau.

By
\cref{itm:p21p3-gf-bridge:ray,itm:p21p3-gf-bridge:separation}, the target
germ there is the isolated ray
\[
 F_{\mathrm{trop}}(x)+(0,\varepsilon_i)v_{b_i},
\]
and \(F_{\mathrm{trop}}\) identifies its source skeleton with that ray.  Hence,
on a sufficiently short relative-open subsegment, the Euclidean current
\[
 \langle T_{\Omega},\beta\rangle
       =\langle T|_U,F^*\beta\rangle
\]
used in the proof of \cite[Theorem~16.8.5]{CLD25} receives the localized
branch contribution from precisely that virtual-annular collar.  Indeed, a
point on a different localized source piece with the same concatenated
tropical value would have the same values under every appended plateau
tropicalization.  The plateau itself is a smooth function of those moment
coordinates, so it would have to equal one there as well, contrary to the
disjoint source localization.  The
primitive scalar slope of the appended norm-one lift of a residue \(g\) is
the corresponding coordinate of the enlarged \(v_{b_i}\), namely
\(\operatorname{ord}_{p_{b_i}}(g)\) by
\eqref{eq:p21p3-gf-branch-order-residue}.  This proves
\cref{itm:p21p3-gf-bridge:localization}.

Equation \eqref{eq:p21p3-gf-common-expansion-index} also proves the claim in
\cref{rem:p21p3-gf-virtual-annulus-index}.  Under pushforward, the rank-one
matrix current has one factor \(e_b\): its tangent tensor contributes
\(e_b^2\) and one-dimensional change of measure contributes \(e_b^{-1}\).
The canonical calibration has the same expansion factor by
\cite[\S14.3.2, Lemma~14.3.4, and Proposition~14.4.2]{CLD25}.  No arbitrary
current was extended to a splitting field at any point of the argument.
\end{proof}
